% Template thesis MSc International Economics - University of Milano-Bicocca. created by C. Angelini and A. Colombo 




% Document formatted as indicated by the University
% (\documentclass lives in main.tex so Overleaf detects main.tex as the main document)
\usepackage[top=2.5cm, bottom=2.5cm, inner=3cm, outer=2.5cm, centering]{geometry}
\usepackage{lmodern}% font Times New Roman
\linespread{1.5}
\usepackage{indentfirst}
\setlength{\parindent}{0.5cm} % Rientro del primo rigo del testo a 1.5cm
\usepackage{scrlayer-scrpage} % stili pagina per il frontespizio
\usepackage{titlesec}
\usepackage[hidelinks]{hyperref}
\titleformat{\chapter}{\fontsize{16pt}{18pt}\bfseries}{\thechapter}{1em}{}
\titleformat{\chapter}[block]
  {\normalfont\LARGE\bfseries}{\thechapter.}{0.5em}{\LARGE}
\titlespacing*{\chapter}{0pt}{-20pt}{25pt}
\titleformat{\section}{\fontsize{16pt}{18pt}\bfseries}{\thesection}{1em}{}
\titleformat{\subsection}{\fontsize{14pt}{16pt}\bfseries}{\thesubsection}{1em}{}
\usepackage[bottom]{footmisc} % footNotes 10pt

% Librerie utili
\usepackage{tikz}
\usetikzlibrary{positioning, fit, arrows.meta, backgrounds, shapes.geometric}
\usepackage{longtable}
\usepackage{lipsum} % Solo per generare del testo di riempimento
\usepackage{float} % per il posizionamento delle immagini
\usepackage{amsmath} % AMS Math Package
\usepackage{amsthm} % Theorem Formatting
\usepackage{subfigure}
\usepackage{booktabs}
\usepackage{caption}
\usepackage{amssymb}
\usepackage[utf8]{inputenc} % codifica UTF-8
\usepackage{xcolor} % Per colorare le righe
\usepackage{colortbl} % Per colorare le tabelle
\usepackage{siunitx}
\usepackage{rotating}
\usepackage{tablefootnote}
\usepackage{adjustbox}
\usepackage{fullpage}
\usepackage{comment}
\usepackage{tcolorbox}
\tcbuselibrary{skins}
\usepackage{booktabs}
\raggedbottom
\definecolor{procblue}{RGB}{43,87,124}
\definecolor{procorange}{RGB}{243,146,0}
\newcommand{\cnode}[1]{\textcolor{procblue}{\texttt{#1}}}            % sampled clone
\newcommand{\chid}[1]{\textcolor{black!45}{\texttt{\textit{[#1 h.]}}}} % hidden clone
\newcommand{\cmat}[1]{\textcolor{procorange}{\texttt{#1}}}
\clearpairofpagestyles
\cfoot[\pagemark]{\pagemark}
\pagestyle{scrplain}
\usepackage{graphicx} % inserimento di immagini

% ---- Stan code listings (pretty) ----
\usepackage{listings}
\definecolor{stanbg}{RGB}{248,249,251}
\definecolor{stankw}{RGB}{12,68,124}    % blocks / control  (blue)
\definecolor{stanty}{RGB}{15,110,86}    % types             (teal)
\definecolor{stanfn}{RGB}{166,72,46}    % distributions/fns (coral)
\definecolor{stancm}{RGB}{124,124,118}  % comments          (grey)
\definecolor{stanrule}{RGB}{223,227,232}
\lstdefinelanguage{Stan}{
  morekeywords={data,parameters,model,transformed,generated,quantities,functions,
    for,in,while,if,else,return,print,reject,target,real,int,vector,array,matrix,
    row_vector,simplex,ordered,positive_ordered,lower,upper},
  morekeywords=[2]{normal,lognormal,exponential,poisson,poisson_log,neg_binomial_2,
    uniform,gamma,beta,dirichlet,bernoulli,categorical,exp,log,fmin,fmax,sum,
    log_sum_exp,to_vector,rep_vector,dot_product,inv,sqrt},
  sensitive=true, morecomment=[l]{//}, morecomment=[s]{/*}{*/}, morestring=[b]",
}
\lstdefinestyle{stan}{
  language=Stan, backgroundcolor=\color{stanbg}, basicstyle=\ttfamily\scriptsize,
  keywordstyle=\color{stankw}\bfseries, keywordstyle=[2]\color{stanfn},
  commentstyle=\color{stancm}\itshape, stringstyle=\color{stanfn},
  numbers=left, numberstyle=\tiny\color{stancm}, numbersep=9pt,
  showstringspaces=false, breaklines=true, breakatwhitespace=true,
  frame=single, framesep=5pt, rulecolor=\color{stanrule},
  xleftmargin=16pt, framexleftmargin=16pt, aboveskip=1em, belowskip=1em,
  tabsize=2, columns=fullflexible, keepspaces=true,
}
\usepackage{csquotes} % per le citazioni "in blocco"
\usepackage[backend=biber, style=authoryear ]{biblatex} % bibliografia con pacchetto biblatex (https://ctan.org/pkg/biblatex?lang=en)
\appto{\bibsetup}{\raggedright}
\addbibresource{bibliography.bib}

% ---- Theorem-like environments (statistics-thesis style) ----
\theoremstyle{plain}
\newtheorem{theorem}{Theorem}[chapter]
\newtheorem{lemma}{Lemma}[chapter]
\newtheorem{corollary}{Corollary}[chapter]
\theoremstyle{definition}
\newtheorem{definition}{Definition}[chapter]
\newtheorem{assumption}{Assumption}[chapter]
\theoremstyle{remark}
\newtheorem{remark}{Remark}[chapter]

